% ---------------------------------------------------------------
% When the Attack Is Not an Attack
% Compiles standalone with pdflatex. For an ACL venue, swap
% \documentclass for \usepackage[review]{acl} and drop geometry+natbib.
%
% TODO before submission: add related work on automated red-teaming
% and LLM-as-evaluator reliability. Only cite papers you have read.
% ---------------------------------------------------------------
\documentclass[11pt]{article}
\usepackage[margin=0.95in]{geometry}
\usepackage{times}
\usepackage{booktabs}
\usepackage{amsmath}
\usepackage{amssymb}
\usepackage{url}
\usepackage[hidelinks]{hyperref}
\usepackage{microtype}
\usepackage{natbib}   % ACL templates load this already; drop when switching

\title{When the Attack Is Not an Attack:\\
Validity Failures in LLM-Generated Red-Teaming}

\author{Yihang Jiao \\
  University of Illinois Urbana-Champaign \\
  \texttt{yihangj3@illinois.edu}}

\date{}

\begin{document}
\maketitle

\begin{abstract}
Language models are increasingly used to generate the probes that test other
models. The appeal is scale and fluency: an LLM can rewrite an adversarial
example into text a person might actually write. This paper asks what such a
pipeline is measuring. I identify three ways a generated probe can be counted in
a red-teaming result without being a valid test, and I quantify each on 32
model--attack--dataset configurations and 600 fresh generations. \textbf{Label
validity}: the rewriting prompt that produces the most label flips is the one
that instructs the model to change the label, so the probe is a different input
rather than an attack. Reporting flip rate alone overstates measured attack
success by $3.05\times$ overall and by $19.5\times$ on topic classification.
\textbf{Input validity}: on sentence-pair tasks most rewrites collapse the two
segments and are not well-formed classifier inputs; because malformed probes are
scored as failures, leaving them in \emph{understates} attack success, 12.5\%
against 41.7\%. The two errors therefore push in opposite directions and do not
cancel. \textbf{Measurement validity}: the reported metric must be a function of
the generated text, a condition an earlier version of this study did not meet.
Underneath these, LLM rewriting is a trade: it keeps only 44.5\% of the attacks
it is applied to while making them far more fluent (median perplexity 22.9
against 522). I close with three numbers any automated red-teaming report should
carry.
\end{abstract}

\section{Introduction}

Automated red-teaming increasingly uses one language model to probe another. The
attraction is obvious. Word-substitution attacks such as
TextFooler~\citep{jin2020bert} and PWWS~\citep{ren2019pwws} produce inputs that
are effective but stilted, and a perturbation no fluent speaker would write says
less about deployment risk than one that reads normally. An LLM can rewrite the
whole sentence rather than swap tokens, and can therefore generate probes that
are fluent by construction.

The generated probe, however, only tells us something about the target model if
it is a valid test of that model. This paper asks how often it is, and what the
answer does to the numbers such pipelines report.

A probe can fail to be a valid test in three distinct ways, and each one is
invisible in the summary statistic that red-teaming reports usually carry:

\begin{itemize}
\item \textbf{Label validity.} The rewrite may have changed the input's true
      label. Then a prediction change is correct behaviour, not a vulnerability.
\item \textbf{Input validity.} The rewrite may not be a well-formed input for
      the task at all --- for a sentence-pair task, not a pair.
\item \textbf{Measurement validity.} The reported number may not be a function
      of the generated text, if the pipeline scores something else.
\end{itemize}

I measure all three. The first two are quantified on 600 generations produced for
this paper under controlled prompt conditions; the third is demonstrated on a
worked instance, an earlier version of this study, where it silently determined
the published result.

Underneath the three failures is a simpler question that the same experiments
answer. Rewriting is applied to attacks that \emph{already succeed}, so its cost
is directly measurable: how many survive being rewritten?

My contributions are:

\begin{itemize}
\item \textbf{Rewriting keeps under half the attacks it is applied to.} Across 32
      configurations, LLaMA-2 rewrites of successful word-substitution attacks
      still flip the classifier 44.5\% of the time, while being far more fluent
      (Section~\ref{sec:cost}).
\item \textbf{Label validity fails in the direction that flatters the attacker.}
      A task-specific prompt roughly doubles retention, but on topic
      classification it does so by instructing a topic change. Flip rate alone
      overstates attack success by $3.05\times$ pooled and $19.5\times$ on AG
      News (Section~\ref{sec:label}).
\item \textbf{Input validity fails in the opposite direction.} Malformed
      sentence-pair probes are scored as failures, so not checking them
      understates attack success --- 12.5\% against 41.7\%. An explicit structure
      constraint nearly doubles validity but does not fix it
      (Section~\ref{sec:input}).
\item \textbf{I measure the generate-and-filter pipeline} proposed for this
      setting and report what each constraint removes; the similarity constraint,
      not fluency, does most of the work (Section~\ref{sec:filter}).
\item \textbf{A worked instance of measurement-validity failure}, and the
      one-line check that catches it (Section~\ref{sec:measurement}), followed by
      three numbers a red-teaming report should carry
      (Section~\ref{sec:checklist}).
\end{itemize}

\section{Experimental Setup}
\label{sec:setup}

\paragraph{Attacks and targets.}
I use TextFooler~\citep{jin2020bert} and PWWS~\citep{ren2019pwws} as implemented
in TextAttack~\citep{morris2020textattack}, against \texttt{bert-base-uncased}
and \texttt{roberta-base} classifiers fine-tuned on each dataset. Attacks were
run over 500 test examples per configuration (277--500 after dataset
exhaustion), yielding 190--479 successful attacks per configuration to sample
from.

\paragraph{Datasets.}
Single-sentence tasks: IMDb, SST-2, MR (Rotten Tomatoes), AG News, and CoLA.
Sentence-pair tasks: RTE, QNLI, and MRPC. Pair inputs are two segments joined by
a \texttt{<SPLIT>} marker, following TextAttack's convention.

\paragraph{Probe generation.}
Probes are produced by \texttt{Llama-2-7b-chat-hf}~\citep{touvron2023llama} in
half precision, with nucleus sampling ($p=0.9$, $T=0.8$, 160 new tokens),
through the model's chat template. I strip the scaffolding chat models place at
both ends of an answer (``Here is the revised text:'' and trailing explanations)
before scoring; leaving it in would inflate perplexity and depress similarity for
whichever prompt elicits more chattiness, which is the quantity under study.

\paragraph{Measurements.}
For each probe I record whether the target classifier's prediction differs from
its prediction on the original text (\emph{flip}); cosine similarity between
sentence embeddings of the original and the probe, using
\texttt{all-MiniLM-L6-v2}; and perplexity under GPT-2-XL. For pair tasks I also
record whether the probe preserved the two-segment structure and is a
well-formed input at all.

\paragraph{The comparison of interest.}
Probes are generated only for examples on which the word-substitution attack
already succeeded, so the baseline is fixed at 100\% by construction and the
quantity to report is what fraction of those attacks survives rewriting.

\section{Rewriting Keeps Under Half the Attacks}
\label{sec:cost}

I first re-evaluate a stored corpus of rewrites spanning 32 configurations
(2 target architectures $\times$ 2 attack recipes $\times$ 8 datasets), running
each target classifier on the rewritten text.

On the 20 single-sentence configurations, rewrites still flip the classifier on
average 44.5\% of the time. Every one of these was a successful attack before
rewriting, so more than half are destroyed by it. Retention varies enormously:

\begin{center}
\begin{tabular}{lccc}
\toprule
Dataset & Mean & Min & Max \\
\midrule
AG News & 0.907 & 0.706 & 1.000 \\
IMDb    & 0.764 & 0.667 & 0.889 \\
MR      & 0.268 & 0.077 & 0.462 \\
SST-2   & 0.217 & 0.105 & 0.300 \\
CoLA    & 0.069 & 0.000 & 0.133 \\
\bottomrule
\end{tabular}
\end{center}

The spread is a factor of thirteen, and it lines up with an implementation
detail: AG News and IMDb were the two datasets for which a hand-written prompt
template existed, while the others fell through to a generic one. Dataset and
prompt are confounded here, so the next section separates them --- and in doing
so uncovers the first validity failure.

\section{Label Validity}
\label{sec:label}

To break the confound I run \emph{both} prompt variants on \emph{every} dataset.
The generic prompt is the original fallback (``Improve the text fluency while
altering the classification\ldots''). The task-specific prompts follow the shape
of the hand-written ones: for IMDb and AG News these are the originals; for SST-2
I write one in the same shape, so that ``has a task-specific prompt'' becomes a
manipulable variable rather than a property of the dataset. I sample 30
successful TextFooler attacks per cell.

\begin{table}[t]
\centering
\small
\begin{tabular}{llccl}
\toprule
Dataset & Target & Specific & Generic & $p$ \\
\midrule
AG News & BERT    & 20/30 & 10/30 & 0.019 \\
AG News & RoBERTa & 19/30 &  8/30 & 0.009 \\
IMDb    & BERT    & 22/30 &  7/30 & 0.0002 \\
IMDb    & RoBERTa & 28/30 &  6/30 & $<10^{-4}$ \\
SST-2   & BERT    & 12/30 & 13/30 & 1.000 \\
SST-2   & RoBERTa & 18/30 & 14/30 & 0.438 \\
\midrule
\textbf{Pooled} & & \textbf{119/180} & \textbf{58/180} & $1.7\times10^{-10}$ \\
\bottomrule
\end{tabular}
\caption{Probes that still flip the classifier, by prompt variant, with dataset
and target held fixed. Two-sided Fisher exact test per cell.}
\label{tab:prompt}
\end{table}

Pooled over 360 probes, the task-specific prompt roughly doubles the fraction of
attacks that survive: 119/180 against 58/180 ($p = 1.7\times10^{-10}$). The
effect holds in 5 of 6 cells and is significant in 4; the two SST-2 cells show no
reliable difference.

\paragraph{Why the prompt wins matters more than that it wins.}
An adversarial example must also preserve the input's meaning, and therefore its
true label. Splitting by task type shows two very different mechanisms behind the
same doubling:

\begin{center}
\begin{tabular}{lcc}
\toprule
 & Flip rate & Similarity \\
\midrule
\multicolumn{3}{l}{\emph{Sentiment (IMDb, SST-2)}} \\
\quad specific & 0.667 & 0.713 \\
\quad generic  & 0.333 & 0.705 \\
\midrule
\multicolumn{3}{l}{\emph{Topic (AG News)}} \\
\quad specific & 0.650 & \textbf{0.361} \\
\quad generic  & 0.300 & 0.691 \\
\bottomrule
\end{tabular}
\end{center}

On sentiment the specific prompt doubles the flip rate at no cost in similarity,
0.713 against 0.705. On AG News the same doubling arrives with similarity
collapsing from 0.691 to 0.361 --- and the reason is written in the prompt, which
instructs the model to change the topic ``from world to sports''. A text whose
topic has actually changed is not an adversarial example. It is a different input
with a different correct label, and a prediction change on it is the classifier
behaving correctly.

\paragraph{How much this distorts a report.}
Table~\ref{tab:overstate} puts a number on it. If a red-teaming pipeline reports
flip rate alone, and we then ask how many of those flips came from probes that
still mean what the original meant, the reported figure overstates measured
attack success by $3.05\times$ pooled --- and by $19.5\times$ on the task where
the prompt instructs a semantic change.

\begin{table}[t]
\centering
\small
\begin{tabular}{lccc}
\toprule
Dataset & Flip rate & Meaning-preserving & Ratio \\
\midrule
AG News & 65.0\% & \textbf{3.3\%}  & $19.5\times$ \\
SST-2   & 50.0\% & 15.0\%          & $3.3\times$ \\
IMDb    & 83.3\% & 46.7\%          & $1.8\times$ \\
\midrule
\textbf{Pooled} & \textbf{66.1\%} & \textbf{21.7\%} & $\mathbf{3.05\times}$ \\
\bottomrule
\end{tabular}
\caption{Attack success under the task-specific prompt as a flip-rate-only report
would state it, against the same probes required to keep similarity $\geq 0.75$.}
\label{tab:overstate}
\end{table}

This is a trap for the whole approach, and it is structural rather than
accidental: \textbf{the instruction most effective at producing label flips is
the instruction to change the label.} Any search over prompts that optimizes
reported attack success will find it.

\section{Input Validity}
\label{sec:input}

For RTE, QNLI and MRPC the classifier expects two segments. Most probes do not
supply them: the model merges premise and hypothesis into a single passage, often
appending a commentary sentence. In the stored 32-configuration corpus, only 100
of 198 pair-task rewrites preserved the structure; for the two RTE configurations
under PWWS, \emph{zero} of 31 did.

I test whether an explicit constraint fixes this, comparing the generic prompt
against one stating that the input is a pair, that the segments are separated by
\texttt{<SPLIT>}, and that the answer must return exactly two sentences in the
same order.

\begin{table}[t]
\centering
\small
\begin{tabular}{lccc}
\toprule
Prompt & Well-formed & Unchecked & Valid only \\
\midrule
generic    & 0.300 & 12.5\% & \textbf{41.7\%} \\
structured & \textbf{0.575} & 21.7\% & 37.7\% \\
\bottomrule
\end{tabular}
\caption{Sentence-pair tasks, 240 probes. \emph{Well-formed}: fraction preserving
the two-segment structure ($p=2.8\times10^{-5}$ between prompts, Fisher exact).
\emph{Unchecked}: attack success if malformed probes are left in the denominator.
\emph{Valid only}: attack success among well-formed probes.}
\label{tab:pair}
\end{table}

Two things follow. First, stating the constraint nearly doubles structural
validity, from 30.0\% to 57.5\%, and raises the yield of usable probes from
12.5\% to 21.7\% --- but the model still violates an explicitly stated format
more than 40\% of the time, so the constraint is a mitigation and not a fix.

Second, and more consequential for anyone reading a red-teaming number:
malformed probes are counted as attack failures, never as successes --- none of
the malformed probes registered a flip. Leaving them in the denominator therefore
\textbf{understates} attack success, 12.5\% against 41.7\% for the generic
prompt. Label-validity failure inflates the reported number and input-validity
failure deflates it. \textbf{The two errors push in opposite directions, so they
do not cancel}; which one dominates depends on the task, and neither is visible
in the summary statistic.

\section{What the Filter Trades}
\label{sec:filter}

The standard response to unreliable generations is to filter them: keep only
probes that flip the label, stay semantically close, and read fluently. I
elaborate that filter into a per-candidate test --- each probe scored on all
three conditions and kept or dropped individually --- and apply it to all 360
single-sentence probes.

\begin{table}[t]
\centering
\small
\begin{tabular}{lccc}
\toprule
Stage & Kept & Frac. & Sim. / PPL \\
\midrule
All probes               & 360 & 1.000 & 0.648 / 33.8 \\
\;+ flips the classifier & 177 & 0.492 & 0.596 / 28.1 \\
\;+ similarity $\geq$ 0.75 & 64 & 0.178 & 0.840 / 27.2 \\
\;+ perplexity $\leq$ 150  & 55 & 0.153 & 0.834 / 22.9 \\
\bottomrule
\end{tabular}
\caption{Cumulative effect of the three conditions. Sim.\ is mean cosine
similarity, PPL median GPT-2-XL perplexity, of the surviving set.}
\label{tab:filter}
\end{table}

The filter retains 15.3\% of generations and does what it claims: mean similarity
rises from 0.648 to 0.834. Two observations matter for practice. First, the
\emph{similarity} constraint does nearly all the filtering --- it removes 113 of
the 177 probes that flip, while the fluency bound removes 9 more. This inverts
the usual framing, in which fluency is the reason to filter; here fluency is
nearly free and meaning preservation is what costs.

Second, the fluency gain over the attacks being replaced is real and large. Over
428 attacked examples, the median GPT-2-XL perplexity of a TextFooler
perturbation is 522, against 22.9 for filtered rewrites, and only 15.4\% of 403
successful word-substitution attacks satisfy the same similarity and fluency
constraints. \textbf{LLM rewriting buys a large improvement in fluency and pays
for it in attack retention.} Neither side of the trade is small, and the exchange
rate is set by the prompt.

\section{Measurement Validity: A Worked Instance}
\label{sec:measurement}

The third failure is the one that most easily escapes notice, because it leaves no
trace in the outputs. An earlier version of this study reported that filtered
LLaMA-2 generations outperformed word-substitution attacks, with improvements up
to 19.31\% relative. Those numbers do not survive scrutiny, and the reasons are
worth recording because the failure mode is easy to reproduce.

\paragraph{The classifier was never run on the probes.}
The pipeline wrote LLaMA output into a \texttt{llama\_text} column, and the
reported attack success rate was computed from TextAttack's
\texttt{result\_type} field. That field describes the word-substitution attack,
which ran \emph{before} rewriting. The probes therefore never entered the reported
numbers. Section~\ref{sec:cost} supplies the missing evaluation.

\paragraph{The headline table compared sample sizes.}
Each ``filtered'' figure was the attack's success rate over 20 examples; each
``baseline'' figure was the same attack over 277--500. All five rows reproduce
exactly from those two runs, so the reported improvement was a difference in
sample size.

\paragraph{Filtering happened at a different layer than described.}
The pipeline did constrain and score generations, but not where the published
algorithm places it. A similarity constraint, \texttt{min\_cos\_sim} $=0.6$, was
passed to the attack's word-substitution step, and perplexity and grammar were
computed afterwards as aggregate statistics over each result file. A
per-candidate test on all three thresholds is not in that code, and the values
$0.75$ and $150$ appear nowhere in it. Section~\ref{sec:filter} implements the
per-candidate test.

\paragraph{The check that catches it.}
The lesson is narrow and mechanical: \emph{if a pipeline produces new text, the
reported metric must be a function of that text.} A one-line assertion --- that
the evaluated column is the generated column --- would have caught all three
problems. That this happened in a study whose entire subject was generated text
is the point: measurement-validity failure is not a sign of carelessness about
the topic, it is a default of the tooling, which will happily join a generations
column to a metric computed from something else.

\section{Three Numbers a Red-Teaming Report Should Carry}
\label{sec:checklist}

Each failure above is invisible in an attack success rate, and each is visible in
one additional number:

\begin{enumerate}
\item \textbf{Semantic similarity beside every flip rate.} Without it, a prompt
      that changes the label scores best. On AG News this is the difference
      between reporting 65.0\% and 3.3\%.
\item \textbf{The fraction of generations that are well-formed inputs.} Without
      it, malformed probes sit in the denominator as failures. On sentence-pair
      tasks this is the difference between 12.5\% and 41.7\%.
\item \textbf{The number of generations actually scored by the target model.}
      Without it, a pipeline can report a metric computed from something other
      than what it generated, and nothing in the output will look wrong.
\end{enumerate}

\paragraph{Limitations.}
Probes come from a single 7B model; larger or format-tuned models may preserve
structure better. I sample 30 per cell, adequate for the effect sizes reported
but not for small ones, which is why the SST-2 cells are inconclusive rather than
negative. Similarity is measured with one sentence encoder, and encoder choice
shifts how severe the 0.75 threshold is. The task-specific prompt for SST-2 was
written by me in the shape of the originals; a different phrasing might move the
effect. Finally, the three failures are demonstrated on classification; whether
they take the same form in generative red-teaming is untested here.

\section{Conclusion}

An automated red-teaming result is a claim about a target model, but it is
produced by a second model whose output may not be a valid test at all. I
measured three ways that happens. A prompt can raise the flip rate by changing
the input's true label, which overstates attack success by $3.05\times$ pooled
and $19.5\times$ on topic classification. A probe can fail to be a well-formed
input, which understates it, 12.5\% against 41.7\%. And a pipeline can report a
number computed from something other than what it generated, which is
undetectable from the outputs alone.

Underneath, LLM rewriting is a trade rather than an improvement: it makes attacks
dramatically more fluent, median perplexity 22.9 against 522, and destroys more
than half of them. None of this makes the approach unusable. It makes it
measurable, which requires reporting similarity beside flip rate, well-formedness
beside both, and how many generations the target model actually saw.

\section*{Reproducibility}

Code, the stored probes, and every result table in this paper are available at
\url{https://github.com/yhj3/counterfactual-textattack}. The repository documents
which components are original and which were reconstructed, and includes the
audit described in Section~\ref{sec:measurement}.

\begin{thebibliography}{}

\bibitem[Jin et al.(2020)]{jin2020bert}
Di Jin, Zhijing Jin, Joey Tianyi Zhou, and Peter Szolovits. 2020.
\newblock Is BERT really robust? A strong baseline for natural language attack
  on text classification and entailment.
\newblock In \emph{AAAI}, volume 34, pages 8018--8025.

\bibitem[Morris et al.(2020)]{morris2020textattack}
John X. Morris, Eli Lifland, Jin Yong Yoo, Jake Grigsby, Di Jin, and Yanjun Qi.
  2020.
\newblock TextAttack: A framework for adversarial attacks, data augmentation,
  and adversarial training in NLP.
\newblock In \emph{EMNLP: System Demonstrations}, pages 119--126.

\bibitem[Ren et al.(2019)]{ren2019pwws}
Shuhuai Ren, Yihe Deng, Kun He, and Wanxiang Che. 2019.
\newblock Generating natural language adversarial examples through probability
  weighted word saliency.
\newblock In \emph{ACL}, pages 1085--1097.

\bibitem[Touvron et al.(2023)]{touvron2023llama}
Hugo Touvron, Louis Martin, Kevin Stone, et al. 2023.
\newblock Llama 2: Open foundation and fine-tuned chat models.
\newblock \emph{arXiv preprint arXiv:2307.09288}.

\bibitem[Ebrahimi et al.(2018)]{ebrahimi2018hotflip}
Javid Ebrahimi, Anyi Rao, Daniel Lowd, and Dejing Dou. 2018.
\newblock HotFlip: White-box adversarial examples for text classification.
\newblock In \emph{ACL (Short Papers)}, pages 31--36.

\bibitem[Kaushik et al.(2020)]{kaushik2020learning}
Divyansh Kaushik, Eduard Hovy, and Zachary C. Lipton. 2020.
\newblock Learning the difference that makes a difference with
  counterfactually-augmented data.
\newblock In \emph{ICLR}.

\end{thebibliography}

\end{document}
